%% Chapter 3 — Site Context
\chapter{Site Context}
\label{ch:sitecontext}

\section{Location and Extent}

The study site is centred on the property known as Pingle, Taylor's Lane, Ashleyhay, Derbyshire, DE4~4AH. The site centroid falls at approximately National Grid Reference SK~310~525 (WGS84: \dataplaceholder{LAT}°N, \dataplaceholder{LON}°W). The analysis covers a circular study area of 1\,km radius from this centroid, encompassing approximately 314\,ha of the surrounding landscape.

\begin{figure}[H]
\centering
\begin{tikzpicture}
  \draw[dreamlab-light,fill=dreamlab-light,rounded corners] (0,0) rectangle (12,8);
  \node[dreamlab-grey,font=\sffamily\large] at (6,4) {Figure placeholder: Site location map};
  \node[dreamlab-grey,font=\sffamily\small] at (6,3) {OS base mapping with 1\,km study radius overlay};
  \node[dreamlab-grey,font=\sffamily\footnotesize] at (6,2) {../output/figures/site\_location\_map.png};
\end{tikzpicture}
\caption{Location of the Pingle study site within the Derbyshire landscape, showing the 1\,km radius study boundary (red circle), principal access routes, and nearby settlements. Base mapping: \copyright\ Crown Copyright and database rights 2026 Ordnance Survey.}
\label{fig:site-location}
\end{figure}

\subsection{Administrative Context}

The site falls within the following administrative boundaries:

\begin{table}[H]
\centering
\caption{Administrative context for the study area.}
\label{tab:admin-context}
\begin{tabular}{ll}
\toprule
\textbf{Administrative Unit} & \textbf{Value} \\
\midrule
County                & Derbyshire \\
District              & Amber Valley \\
Parish                & Ashleyhay (part), Wirksworth (part) \\
Parliamentary         & Derbyshire Dales \\
Planning Authority    & Amber Valley Borough Council \\
Minerals Authority    & Derbyshire County Council \\
\bottomrule
\end{tabular}
\end{table}


\section{Landscape Character}

The study area lies within the Peak Fringe and Lower Derwent character area, at the transition zone between the White Peak limestone plateau to the north-west and the lower-lying agricultural lands of the Trent Valley to the south-east. The landscape is characterised by:

\begin{itemize}[nosep]
  \item Rolling pastoral farmland with a mix of improved and semi-improved grasslands;
  \item Scattered broadleaved woodland, including ancient woodland fragments;
  \item A network of hedgerows, many species-rich and of historical significance;
  \item Small streams and seasonal watercourses draining south-east towards the Ecclesbourne valley;
  \item Dispersed settlement pattern with farmsteads and isolated dwellings connected by narrow lanes.
\end{itemize}

The Derbyshire Landscape Character Assessment (2014) classifies the area as Landscape Character Type~7: Settled Farmlands, Sub-type~7b: Slopes and Valleys with Woodland. Key characteristics include the small-scale field pattern enclosed by hedgerows, the presence of veteran trees along field boundaries, and the intimate, enclosed landscape character.


\section{Statutory and Non-Statutory Designations}

\subsection{Designations Within the Study Area}

\begin{table}[H]
\centering
\caption{Environmental designations intersecting or proximate to the 1\,km study area.}
\label{tab:designations}
\begin{tabular}{p{40mm}p{30mm}p{60mm}}
\toprule
\textbf{Designation} & \textbf{Status} & \textbf{Notes} \\
\midrule
Derwent Valley Mills World Heritage Site & UNESCO WHS  & Buffer zone extends to within \dataplaceholder{DIST}\,m of study area boundary \\
Gang Mine SSSI            & Statutory         & Located \dataplaceholder{DIST}\,km to the north-west; lead mine with calaminarian grassland \\
Alport Height LWS         & Non-statutory     & Local Wildlife Site adjacent to study area boundary \\
Wirksworth SSSI           & Statutory         & Limestone quarry SSSI \dataplaceholder{DIST}\,km south \\
Ancient Woodland Inventory & Non-statutory    & \dataplaceholder{N} parcels within study area identified as ancient or semi-natural woodland \\
\bottomrule
\end{tabular}
\end{table}

\begin{confidencebox}{HIGH}
Designation boundaries were obtained from Natural England Open Data (Magic Map portal, accessed March 2026) and cross-referenced with the Amber Valley Local Plan Proposals Map.
\end{confidencebox}

\subsection{Local Nature Recovery Strategy}

Derbyshire County Council's emerging Local Nature Recovery Strategy (LNRS), currently at draft v0.3 (December 2025), identifies habitat creation and enhancement priorities across the county. Within the study area, the LNRS draft identifies:

\begin{itemize}[nosep]
  \item \dataplaceholder{N} parcels within a grassland enhancement priority zone;
  \item \dataplaceholder{N} parcels within a woodland connectivity corridor;
  \item The Ecclesbourne tributary as a priority watercourse for riparian habitat restoration.
\end{itemize}

\begin{confidencebox}{LOW}
The LNRS is not yet adopted. Strategic significance assignments based on the draft v0.3 layer may change following public consultation and adoption, currently expected Q3 2026.
\end{confidencebox}


\section{Geology and Soils}

\subsection{Bedrock Geology}

The study area spans the geological transition from Carboniferous Limestone (Bee Low Limestone Formation) in the north-west to Millstone Grit Series (Ashover Grit) in the south-east. The British Geological Survey 1:50,000 mapping (Sheet~112, Chesterfield) shows:

\begin{itemize}[nosep]
  \item \textbf{North-western quadrant}: Bee Low Limestone Formation — pale grey, massively bedded limestone with chert bands. Supports calcareous grassland where soils are thin.
  \item \textbf{Central band}: Eyam Limestone Formation — dark grey, thinly bedded limestone with shale partings.
  \item \textbf{South-eastern quadrant}: Ashover Grit Formation — coarse-grained sandstone and siltstone. Supports more acidic soils and heathland/acid grassland habitats.
\end{itemize}

\subsection{Superficial Deposits}

Superficial deposits are patchy across the study area, consisting primarily of:
\begin{itemize}[nosep]
  \item Head deposits (clay, silt, and gravel) in valley bottoms;
  \item Glaciofluvial sand and gravel deposits on lower slopes;
  \item Thin or absent cover on hilltops and upper slopes, where bedrock is near-surface.
\end{itemize}

\subsection{Soils}

The National Soil Map (Cranfield University, 1:250,000) classifies the soils within the study area as:
\begin{itemize}[nosep]
  \item \textbf{Aberford association} (north-west): shallow, well-drained calcareous soils over limestone. Brown rendzinas and calcareous brown earths.
  \item \textbf{Rivington~2 association} (south-east): well-drained coarse loamy soils over sandstone. Typical brown earths.
  \item \textbf{Cegin association} (valleys): slowly permeable seasonally waterlogged soils. Stagnogleys supporting rush pasture.
\end{itemize}

The soil association strongly influences the habitat types present, with the limestone soils supporting species-rich calcareous grassland and the sandstone soils supporting acid grassland and heathland communities.


\section{Hydrology}

The study area is drained by a small unnamed tributary of the Ecclesbourne, flowing south-east through the central valley. The watercourse is classified as an ordinary watercourse under the jurisdiction of Derbyshire County Council as Lead Local Flood Authority. Key hydrological features include:

\begin{itemize}[nosep]
  \item A perennial stream approximately \dataplaceholder{LENGTH}\,m in length within the study boundary;
  \item Two seasonal drainage channels on the northern slopes;
  \item A small farm pond (approximately \dataplaceholder{AREA}\,m\textsuperscript{2}) in the south-western quadrant;
  \item No Environment Agency main river designation within the study area.
\end{itemize}

The watercourse and its immediate riparian zone are assessed as linear features under the BM\,4.0 watercourse module (Chapter~\ref{ch:linear}).

\begin{figure}[H]
\centering
\begin{tikzpicture}
  \draw[dreamlab-light,fill=dreamlab-light,rounded corners] (0,0) rectangle (12,8);
  \node[dreamlab-grey,font=\sffamily\large] at (6,4) {Figure placeholder: Hydrology and drainage map};
  \node[dreamlab-grey,font=\sffamily\footnotesize] at (6,3) {../output/figures/hydrology\_map.png};
\end{tikzpicture}
\caption{Hydrology and drainage features within the 1\,km study area, showing the perennial watercourse (blue), seasonal channels (dashed blue), and the farm pond (circle).}
\label{fig:hydrology}
\end{figure}


\section{Planning History}

A review of the Amber Valley Borough Council planning register for the study area reveals the following relevant planning history:

\begin{longtable}{p{25mm}p{20mm}p{75mm}}
\caption{Relevant planning history within the study area.}
\label{tab:planning-history} \\
\toprule
\textbf{Application Ref} & \textbf{Date} & \textbf{Description \& Decision} \\
\midrule
\endfirsthead
\midrule
\endhead
\bottomrule
\endlastfoot

\dataplaceholder{REF} & \dataplaceholder{DATE} & Agricultural building — Approved. No habitat conditions. \\
\dataplaceholder{REF} & \dataplaceholder{DATE} & Residential extension — Approved. Standard conditions. \\
\dataplaceholder{REF} & \dataplaceholder{DATE} & Change of use agricultural to equestrian — Approved with landscape management plan condition. \\
\dataplaceholder{REF} & \dataplaceholder{DATE} & Outline application for residential development — Refused. Ecological impact cited as a reason for refusal. \\
\end{longtable}

\begin{confidencebox}{MEDIUM}
Planning history was obtained from the Amber Valley Borough Council online planning register, accessed March 2026. Historic applications prior to 2005 may not be available online and are not included.
\end{confidencebox}


\section{Land Use Summary}

Based on OS MasterMap descriptive terms and initial aerial photograph interpretation, the study area comprises the following broad land use categories:

\begin{table}[H]
\centering
\caption{Broad land use composition within the 1\,km study area.}
\label{tab:landuse-summary}
\begin{tabular}{lrr}
\toprule
\textbf{Land Use Category} & \textbf{Area (ha)} & \textbf{\% of Study Area} \\
\midrule
Grassland (improved + semi-improved) & \dataplaceholder{AREA} & \dataplaceholder{PCT} \\
Woodland (broadleaved + mixed)       & \dataplaceholder{AREA} & \dataplaceholder{PCT} \\
Arable                               & \dataplaceholder{AREA} & \dataplaceholder{PCT} \\
Hedgerow network (area equivalent)   & \dataplaceholder{AREA} & \dataplaceholder{PCT} \\
Built-up / gardens                   & \dataplaceholder{AREA} & \dataplaceholder{PCT} \\
Water features                       & \dataplaceholder{AREA} & \dataplaceholder{PCT} \\
Other (tracks, bare ground)          & \dataplaceholder{AREA} & \dataplaceholder{PCT} \\
\midrule
\textbf{Total}                       & \dataplaceholder{TOTAL} & 100.0 \\
\bottomrule
\end{tabular}
\end{table}
